% ============================================================
%  Apéndice B — Archivos de salida
%  TT 2026-B066 | ESCOM-IPN
%  Fuente: inventario del sistema (pipeline/, worker/, backend-spring/)
%  Los ejemplos son de forma: nombres y tipos reales, valores ilustrativos.
% ============================================================

\chapter{Archivos de salida del análisis}
\label{ap:salidas}

Este apéndice describe los archivos que produce el análisis de un video y los que el
investigador obtiene del sistema. No hay un único archivo de seguimiento: la etapa de
cilindros y el clasificador escriben tablas CSV de trabajo, el servicio de análisis
escribe unos pocos archivos JSON de apoyo y los resultados definitivos se guardan en
la base de datos. Los reportes que descarga el investigador se arman desde la base, no
desde los CSV de trabajo.

Los ejemplos de este apéndice muestran la forma de cada archivo: los nombres, el orden
y el tipo de las columnas son los reales, pero los valores son ilustrativos.

\section{Dónde se guardan}

Todo lo de un análisis vive en el volumen \texttt{videos\_reportes}, en la carpeta
\texttt{analisis/\{id del análisis\}/}. Los archivos de la etapa de cilindros y del
clasificador van en su subcarpeta \texttt{trabajo/} y su nombre empieza con el
identificador del video.

\begin{longtable}{|p{4.3cm}|p{6.2cm}|p{3.3cm}|}
\caption{Archivos que genera el análisis.}
\label{tab:ap_archivos}\\
\hline
\textbf{Archivo} & \textbf{Contenido} & \textbf{¿Lo ve el investigador?} \\
\hline
\endfirsthead
\caption[]{Archivos que genera el análisis (continuación).}\\
\hline
\textbf{Archivo} & \textbf{Contenido} & \textbf{¿Lo ve el investigador?} \\
\hline
\endhead
\texttt{\_cuadros.csv} & Medidas de cada fotograma y de cada cilindro. & No \\
\hline
\texttt{\_rasgos.csv} & Los 15 rasgos por bloque de 5~s y por espécimen. & No \\
\hline
\texttt{\_cnnseg.csv} & La conducta de cada segundo, salida de la red. & Solo en modo
desarrollador \\
\hline
\texttt{\_cnnseg\_resumen.csv} & Segundos de cada conducta por minuto. & Solo en modo
desarrollador \\
\hline
\texttt{cajas.json} & Recuadro de cada cilindro, en píxeles del video. & Su efecto:
los recuadros de la revisión \\
\hline
\texttt{web.mp4} & Copia del video, sin audio, de los últimos 300~s. & Sí: es el video
que se reproduce en la revisión \\
\hline
\texttt{progreso.json} & Etapa y porcentaje de avance. & Sí, como barra de avance \\
\hline
\texttt{pausa.json}, \texttt{pausa.jpg} & Estado de la pausa del análisis. & Sí, como
pantalla de pausa \\
\hline
\texttt{correcciones.json} & Ajustes de recuadros hechos en las pausas. & No \\
\hline
\texttt{tiempos.json} & Duración de cada etapa. & Solo en modo desarrollador \\
\hline
Reportes PDF, CSV y XLSX; Excel del experimento & Resultados armados desde la
base. & Sí \\
\hline
\end{longtable}

\section{Archivos de trabajo del clasificador}

\paragraph{\texttt{\_cuadros.csv}.} Una fila por fotograma medido y por cilindro. Las
medidas están en píxeles del video; la normalización se hace después.

\begin{longtable}{|p{2.8cm}|p{2.2cm}|p{8.8cm}|}
\caption{Columnas de \texttt{\_cuadros.csv}.}
\label{tab:ap_cuadros}\\
\hline
\textbf{Campo} & \textbf{Unidad} & \textbf{Significado} \\
\hline
\endfirsthead
\caption[]{Columnas de \texttt{\_cuadros.csv} (continuación).}\\
\hline
\textbf{Campo} & \textbf{Unidad} & \textbf{Significado} \\
\hline
\endhead
\texttt{paso} & Fotogramas & Cada cuántos fotogramas se midió (12 por omisión). \\
\hline
\texttt{t} & Segundos & Instante del fotograma desde el inicio del video. \\
\hline
\texttt{s} & Sin unidad & Número del cilindro, que es el del espécimen. \\
\hline
\texttt{cx}, \texttt{cy} & Píxeles & Centro del espécimen; vacío si no se vio. \\
\hline
\texttt{area} & Píxeles$^2$ & Área de la silueta completa. \\
\hline
\texttt{above} & Píxeles$^2$ & Silueta por encima de la línea de agua. \\
\hline
\texttt{largo} & Píxeles & Largo del cuerpo en ese fotograma. \\
\hline
\texttt{vert} & 0 a 1 & Qué tan vertical está el eje mayor del cuerpo. \\
\hline
\texttt{elong} & 1 o más & Qué tan alargado está el cuerpo. \\
\hline
\texttt{cabeza} & Píxeles & Fila del punto más alto de la silueta sin cola. \\
\hline
\texttt{me} & Niveles de gris & Cambio de imagen dentro del espécimen respecto a
unos 0.4~s antes. \\
\hline
\texttt{iou}, \texttt{dice} & 0 a 1 & Solape de la silueta con la de unos 0.4~s
antes. \\
\hline
\texttt{gx0}, \texttt{gx1} & Píxeles & Límites izquierdo y derecho del cilindro. \\
\hline
\texttt{g\_agua}, \texttt{g\_fondo} & Píxeles & Fila de la línea de agua y del
fondo. \\
\hline
\texttt{escala} & Píxeles & Mediana del largo del cuerpo en todo el video; es la
unidad de longitud. \\
\hline
\end{longtable}

Cuando no hay dato, por ejemplo en las medidas de movimiento de un fotograma sin
espécimen, el campo queda vacío.

\paragraph{\texttt{\_rasgos.csv}.} Una fila por bloque de 5~s y por espécimen. Además
de los 15 rasgos de la sección~\ref{sec:medidas_red}, trae el video, el número de
bloque, el espécimen, el inicio y el fin del bloque en segundos, los fotogramas
medidos y los fotogramas en que se vio al espécimen (con menos de 5 el bloque queda
sin rasgos), la escala en píxeles, la geometría detectada del cilindro y el recuadro
del cilindro en el video original.

\paragraph{\texttt{\_cnnseg.csv}.} Es la decisión del clasificador: una fila por
espécimen y por segundo de todo el video. De aquí el servicio toma los últimos 300
segundos de cada espécimen.

\begin{table}[H]
\centering
\caption{Columnas de \texttt{\_cnnseg.csv}.}
\label{tab:ap_cnnseg}
\begin{tabular}{|p{3.4cm}|p{2.2cm}|p{8.2cm}|}
\hline
\textbf{Campo} & \textbf{Unidad} & \textbf{Significado} \\
\hline
\texttt{especimen} & Sin unidad & Número del cilindro. \\
\hline
\texttt{segundo} & Segundos & Segundo del video, desde 1. \\
\hline
\texttt{inicio\_s}, \texttt{tiempo} & Segundos, min:seg & Dónde empieza ese segundo,
en dos formatos. \\
\hline
\texttt{conducta} & Texto & \texttt{escalamiento}, \texttt{inmovilidad},
\texttt{nado} o \texttt{activa}, según la regla de tres pasos con umbral de 0.80. \\
\hline
\texttt{confianza} & 0 a 1 & Probabilidad de la conducta elegida. \\
\hline
\texttt{p\_escalamiento}, \texttt{p\_inmovilidad}, \texttt{p\_nado} & 0 a 1 &
Probabilidad de cada conducta. \\
\hline
\texttt{ventana\_s} & Segundos & Ventana de los rasgos (4.0). \\
\hline
\end{tabular}
\end{table}

Ejemplo de forma:

{ootnotesize
\begin{verbatim}
especimen,segundo,inicio_s,tiempo,conducta,confianza,p_escalamiento,p_inmovilidad,p_nado,ventana_s
1,1,0,0:00,inmovilidad,0.912,0.021,0.912,0.067,4.0
1,2,1,0:01,activa,0.874,0.402,0.126,0.472,4.0
\end{verbatim}}

Al guardar el resultado en la base, \texttt{nado} pasa a «nado activo» y \texttt{activa}
a «conducta activa». Un segundo en el que no se vio al espécimen queda sin conducta,
aunque la red haya dado una.

\paragraph{\texttt{\_cnnseg\_resumen.csv}.} Una fila por espécimen y por minuto, y una
fila \texttt{total} por espécimen, con los segundos de cada conducta
(\texttt{nado}, \texttt{inmovilidad}, \texttt{escalamiento}, \texttt{activa}) y los
segundos sin decidir (\texttt{sin\_decidir}).

\section{Archivos JSON de apoyo}

\paragraph{\texttt{cajas.json}.} Lo escribe el servicio de análisis al completar el
análisis. El servidor lo lee para dibujar los recuadros sobre el video en la revisión;
si falta, usa un tamaño por omisión y deja los recuadros vacíos, sin fallar.

\begin{table}[H]
\centering
\caption{Campos de \texttt{cajas.json}.}
\label{tab:ap_cajas}
\begin{tabular}{|p{4.6cm}|p{2.2cm}|p{7.0cm}|}
\hline
\textbf{Campo} & \textbf{Unidad} & \textbf{Significado} \\
\hline
\texttt{version} & Sin unidad & Versión del formato. \\
\hline
\texttt{ancho}, \texttt{alto} & Píxeles & Tamaño del video. \\
\hline
\texttt{recuadros} & Lista & Un elemento por espécimen. \\
\hline
\texttt{recuadros[].cilindro} & Sin unidad & Número del cilindro. \\
\hline
\texttt{recuadros[].x0}, \texttt{y0}, \texttt{x1}, \texttt{y1} & Píxeles & Esquinas
superior izquierda e inferior derecha. \\
\hline
\end{tabular}
\end{table}

Ejemplo de forma:

\begin{verbatim}
{
  "version": 1, "ancho": 1920, "alto": 1080,
  "recuadros": [
    {"cilindro": 1, "x0": 150, "y0": 340, "x1": 276, "y1": 912},
    {"cilindro": 2, "x0": 530, "y0": 340, "x1": 659, "y1": 912}
  ]
}
\end{verbatim}

\paragraph{Estado del análisis.} Tres archivos se ven en la página como pantallas y no
se descargan:
\begin{itemize}
\item \texttt{progreso.json}: la etapa (\texttt{preprocesamiento}, \texttt{cilindros} o
\texttt{clasificacion}), el porcentaje de la etapa, el porcentaje general y la hora de
la última actualización.
\item \texttt{pausa.json} y \texttt{pausa.jpg}: el cilindro que se perdió, desde qué
segundo, la imagen de la pausa, la hora en que empezó, los recuadros vigentes y el
ajuste que haga el investigador. Se borran al resolverse la pausa.
\item \texttt{correcciones.json}: la lista de ajustes de recuadro hechos en las pausas,
cada uno con el cilindro, el segundo desde el que vale y las cuatro esquinas.
\end{itemize}

\section{Lo que descarga el investigador}

El investigador obtiene sus resultados desde el experimento, en cuatro formatos que el
servidor arma desde la base de datos:

\begin{enumerate}
\item \textbf{PDF de resultados} del experimento, que se muestra en el navegador.
\item \textbf{CSV} de un grupo o de una tanda.
\item \textbf{XLSX} de un grupo o de una tanda, con el mismo contenido en una hoja
llamada «Resultados».
\item \textbf{Excel del experimento}, con una hoja «Resultados» y una hoja «Segundo a
segundo».
\end{enumerate}

El CSV y el XLSX tienen una fila por espécimen y por minuto, con las columnas de la
tabla~\ref{tab:ap_reporte_csv}. Debajo de las filas traen las filas de resumen
\textit{Media}, \textit{Desviación estándar}, \textit{Varianza} y \textit{Revisión
humana}. Si el grupo mezcla análisis de clasificación precisa y agrupada, agregan un
aviso. El CSV lleva codificación UTF-8 con marca de orden de bytes, y neutraliza los
textos que empiezan con \texttt{=}, \texttt{+}, \texttt{-} o \texttt{@} para que una
hoja de cálculo no los interprete como fórmulas.

\begin{table}[H]
\centering
\caption{Columnas del reporte CSV y XLSX.}
\label{tab:ap_reporte_csv}
\begin{tabular}{|p{3.6cm}|p{2.2cm}|p{8.0cm}|}
\hline
\textbf{Columna} & \textbf{Unidad} & \textbf{Significado} \\
\hline
Grupo & Texto & Etiqueta del grupo. \\
\hline
Tanda & Texto & Letra de la tanda. \\
\hline
Espécimen & Texto & «Espécimen» seguido del número del cilindro. \\
\hline
Tiempo (min:seg) & min:seg & Tramo de un minuto, por ejemplo 0:00-1:00. \\
\hline
Nado & Segundos & Segundos de nado activo en ese minuto. \\
\hline
Escalamiento & Segundos & Segundos de escalamiento. \\
\hline
Inmovilidad & Segundos & Segundos de inmovilidad. \\
\hline
Conducta activa & Segundos & Segundos de conducta activa, cuando no se distinguió
entre nado y escalamiento. \\
\hline
Sin ver & Segundos & Segundos sin conducta: no se vio al espécimen o la revisión lo
marcó como no visible. \\
\hline
\end{tabular}
\end{table}

La hoja «Segundo a segundo» del Excel tiene una fila por espécimen y por segundo, con
las columnas Grupo, Tanda, Espécimen, Tiempo (s), Conducta, Nado activo, Inmovilidad,
Escalamiento, Conducta activa, Sin ver y Origen. Las cinco columnas de conducta valen 1
o 0. El origen dice de dónde salió el dato: \texttt{maquina}, \texttt{confirmada},
\texttt{corregida} o \texttt{dudosa}.

El CSV crudo de la red (\texttt{\_cnnseg.csv}, con probabilidades) solo se descarga con
\texttt{MODO\_DESARROLLADOR=true}. Es la salida sin las correcciones de la revisión y no
los resultados del sistema.

\section{Cómo queda cada segundo en la base de datos}

Los resultados no se quedan en los CSV: el servicio de análisis los inserta en la tabla
\texttt{segundo}, una fila por espécimen y por segundo, y de ella se calculan todos los
reportes. Los totales por minuto no se guardan, se calculan al pedirlos. Solo se
guardan los últimos 300 segundos de cada espécimen, numerados de 1 a 300.

\begin{table}[H]
\centering
\caption{Campos de la tabla \texttt{segundo}.}
\label{tab:ap_segundo}
\begin{tabular}{|p{3.4cm}|p{2.2cm}|p{8.2cm}|}
\hline
\textbf{Campo} & \textbf{Unidad} & \textbf{Significado} \\
\hline
\texttt{idobservacion} & Sin unidad & La observación: un espécimen en un video. \\
\hline
\texttt{segundo} & Segundos & De 1 a 300. \\
\hline
\texttt{conducta} & Texto o nulo & \texttt{nado activo}, \texttt{inmovilidad},
\texttt{escalamiento} o \texttt{conducta activa}. Nulo si no se vio al espécimen o la
revisión lo marcó como no visible. \\
\hline
\texttt{origen} & Texto & \texttt{maquina} (por omisión), \texttt{confirmada},
\texttt{corregida} o \texttt{dudosa}. \\
\hline
\end{tabular}
\end{table}
